\documentclass[10pt]{article}
\usepackage[T1]{fontenc}
\usepackage[utf8]{inputenc}
\usepackage[french]{babel}
\usepackage{amsmath,amssymb,mathrsfs,mathpazo}
\usepackage{geometry}
\geometry{a4paper,margin=2cm}
\usepackage{graphicx,array,tabularx,tabularray,colortbl,multirow,multicol,booktabs,enumitem,pifont,ulem,fancybox,diagbox,cancel}
\usepackage{pstricks,pstricks-add,pst-plot,pst-tree,pst-node,pst-text,pst-func,pst-math,pst-bspline,pst-3dplot,pst-eucl,pst-solides3d}
\let\pscalculate\undefined
\usepackage{pst-bezier}
\usepackage{tikz,pgfplots,tkz-tab}
\usetikzlibrary{arrows,arrows.meta,patterns,calc,positioning,trees,decorations.markings,shapes.geometric,shapes.misc,shapes.arrows,intersections,decorations.pathmorphing,decorations.pathreplacing,backgrounds,shadows,babel}
\pgfplotsset{compat=1.18}
\usepgfplotslibrary{fillbetween}
\usepackage[np]{numprint}
\usepackage{listings,makecell,pas-tableur,twoopt}
\let\barre\undefined
\usepackage[pstricks]{tablvar}
\definecolor{darkgreen}{RGB}{0,100,0}
\providecommand{\red}{\color{red}}
\providecommand{\blue}{\color{blue}}
\providecommand{\green}{\color{green}}
\providecommand{\black}{\color{black}}
\providecommand{\white}{\color{white}}
\providecommand{\gray}{\color{gray}}
\providecommand{\cyan}{\color{cyan}}
\providecommand{\magenta}{\color{magenta}}
\providecommand{\yellow}{\color{yellow}}
\providecommand{\violet}{\color{violet}}
\providecommand{\orange}{\color{orange}}
\providecommand{\euro}{\texteuro}
\providecommand{\EUR}{\texteuro}
\providecommand{\ieme}{\textsuperscript{e}}
\providecommand{\degres}{\ensuremath{^\circ}}
\providecommand{\degree}{\ensuremath{^\circ}}
\providecommand{\vv}[1]{\overrightarrow{#1}}
\providecommand{\Coord}[1]{#1}
\providecommand{\decosix}{}
\providecommand{\starredbullet}{\ensuremath{\bullet}}
\pagestyle{empty}
\setlength{\parindent}{0pt}
\usepackage{adjustbox}\DeclareUnicodeCharacter{2236}{:}\DeclareUnicodeCharacter{2190}{\ensuremath{\leftarrow}}
\providecommand{\R}{\mathbb{R}}
\providecommand{\N}{\mathbb{N}}
\providecommand{\Z}{\mathbb{Z}}
\providecommand{\Q}{\mathbb{Q}}
\providecommand{\D}{\mathbb{D}}
\providecommand{\C}{\mathbb{C}}
\providecommand{\e}{{\mathrm{e}}}
\providecommand{\vect}[1]{\overrightarrow{#1}}
\providecommand{\vectt}[1]{\overrightarrow{\mathrm{#1}}}
\providecommand{\vvt}[1]{\overrightarrow{\mathrm{#1}}}
\providecommand{\np}[2][]{{#2~#1}}
\providecommand{\vv}[1]{{\overrightarrow{#1}}}
\renewcommand{\d}{{\mathrm{d}}}
\providecommand{\degres}{{^\circ}}
\providecommand{\degre}{{^\circ}}
\providecommand{\blue}{\color{blue}}
\providecommand{\red}{\color{red}}
\providecommand{\enskip}{\quad}
\providecommand{\llbracket}{\lbrack\!\lbrack}
\providecommand{\rrbracket}{\rbrack\!\rbrack}
\providecommand{\euro}{{\text{€}}}
\providecommand{\no}{{\text{n°}}}
\providecommand{\interval}[2]{{[#1;#2]}}
\providecommand{\intervFO}[2]{{[#1;#2[}}
\providecommand{\dotfill}{\ldots}

\providecommand{\ds}{\displaystyle}
\providecommand{\textding}[1]{#1}
\providecommand{\Oij}{\(\mathrm{O};\vec{\imath},\vec{\jmath}\)}
\providecommand{\Oijk}{\(\mathrm{O};\vec{\imath},\vec{\jmath},\vec{k}\)}
\providecommand{\Ouv}{\(\mathrm{O};\vec{u},\vec{v}\)}
\providecommand{\intervFF}[2]{\ensuremath{[#1;#2]}}
\renewcommand{\euro}{\ifmmode\text{\texteuro}\else\texteuro\fi}
\DeclareUnicodeCharacter{00B0}{\ifmmode^{\circ}\else\textdegree\fi}
\DeclareUnicodeCharacter{20AC}{\ifmmode\text{\texteuro}\else\texteuro\fi}
% Énoncés : Éducation nationale. Transcription et corrigés : APMEP, Denis Vergès et contributeurs. Découpage et classement : CoopMaths, Sébastien Lozano, Éric Elter et contributeurs. Adaptation de la présentation : Maths Explorer.
% CC BY-SA : https://coopmaths.fr/www/about/
% https://coopmaths.fr/alea/static/eam/2026/tex/eam_technologique_2026_09_metropole_2_cor.tex
\begin{document}
\noindent\textit{APMEP / CoopMaths. CC BY-SA. Présentation adaptée par Maths Explorer.}\par\medskip
\subsection*{Partie A}

\emph{Pour tout réel $x$, $f(x) = (12x - 480)(x - 10)$.}

$12x - 480 = 0 \iff x = \dfrac{480}{12} = 40$ \quad et \quad $x - 10 = 0 \iff x = 10$.

Les deux coefficients de $x$ (12 et 1) sont positifs : chaque
facteur est négatif avant sa racine et positif après.

\begin{center}
	\begin{tikzpicture}
		\tkzTabInit[lgt=3.4,espcl=2.6]
		{$x$/1,$12x-480$/1,$x-10$/1,$f(x)$/1}{$-\infty$,$10$,$40$,$+\infty$}
		\tkzTabLine{,-,t,-,z,+,}
		\tkzTabLine{,-,z,+,t,+,}
		\tkzTabLine{,+,z,-,z,+,}
	\end{tikzpicture}
\end{center}


\subsection*{Partie B}

\emph{Pour $x \in [0 ; 18]$, $B(x) = 4x^3 - 300x^2 + \np{4800}x - \np{10000}$, où $x$ est le nombre de
tracteurs vendus par mois et $B(x)$ le bénéfice en centaines d'euros.}

\begin{enumerate}
	\item $B(0) = 4 \times 0 - 300 \times 0 + \np{4800} \times 0 - \np{10000} = -\np{10000}$.

	Si l'usine ne vend aucun tracteur dans le mois, elle subit une perte
	de \np{10000} centaines d'euros, soit \np{1000000}~\euro.

	\item Lorsque l'usine vend 18 tracteurs dans le mois, elle réalise un bénéfice de \np{2528}
	centaines d'euros, soit \np{252800}~\euro.

	\item
	\begin{enumerate}
		\item Pour tout réel $x$ de $[0~;~18]$ :
		\[B'(x) = 4 \times 3x^2 - 300 \times 2x + \np{4800} = 12x^2 - 600x + \np{4800}.\]
		D'autre part :
		\[f(x) = (12x - 480)(x - 10) = 12x^2 - 120x - 480x + \np{4800} = 12x^2 - 600x + \np{4800}.\]

Donc, pour tout réel $x$ de $[0 ; 18]$, $B'(x) = f(x)$.

		\item Le signe de $B'(x) = f(x)$ se lit dans le tableau de la partie A, restreint à
		l'intervalle $[0 ; 18]$ : seule la racine 10 s'y trouve ($40 > 18$).

		$B(10) = \np{4000} - \np{30000} + \np{48000} - \np{10000} = \np{12000}$.

		\begin{center}
			\begin{tikzpicture}
				\tkzTabInit[lgt=3.4,espcl=2.8]
				{$x$/1,$B'(x)$/1,$B$/2}{$0$,$10$,$18$}
				\tkzTabLine{,+,z,-,}
				\tkzTabVar{-/$-\np{10000}$,+/$\np{12000}$,-/$\np{2528}$}
			\end{tikzpicture}
		\end{center}

		\item D'après le tableau de variations, $B$ atteint son maximum pour $x = 10$ : l'entreprise
		doit produire et vendre \textbf{10 tracteurs} par mois. Le bénéfice maximal est alors
		$B(10) = \np{12000}$ centaines d'euros, soit \np{1200000}~\euro.
	\end{enumerate}
\end{enumerate}

\medskip

\end{document}
